\section{The finite hypergeometric engine}\label{sec:hypergeom}

The residue formulas reduce most of the theory to generalized hypergeometric
series.  We therefore prove their finite theory before mentioning contours.

\subsection{Coefficients and primitive recursion}

For complex parameters $\mathbf A=(A_1,\ldots,A_r)$ and
$\mathbf B=(B_1,\ldots,B_s)$ define
\begin{equation}\label{eq:hypergeom}
 {}_rF_s\!\left(\left.\begin{matrix}\mathbf A\\\mathbf B\end{matrix}\right|w\right)
 =\sum_{k=0}^{\infty}t_k,
 \qquad
 t_0=1,
 \qquad
 t_{k+1}=t_k\,w\,
 \frac{\prod_{i=1}^{r}(A_i+k)}{(k+1)\prod_{j=1}^{s}(B_j+k)}.
\end{equation}
For rational complex inputs, the partial sum and the pair
$(t_k,\sum_{j<k}t_j)$ are obtained by simultaneous primitive recursion.
Denominators are protected by a certificate that no $B_j$ is a nonpositive
integer, or by termination before the first zero when the series is a
polynomial.

\begin{lemma}[Rising-factorial form]\label{lem:rising}
For every $k$ before a denominator zero,
\[
 t_k=\frac{\prod_{i=1}^{r}\poch{A_i}{k}}
           {\prod_{j=1}^{s}\poch{B_j}{k}}\frac{w^k}{k!}.
\]
\end{lemma}
\begin{proof}
The assertion at $k=0$ is a finite ring identity.  Multiplication by the ratio
in \eqref{eq:hypergeom} gives the successor case, using
$\poch A{k+1}=\poch Ak(A+k)$ and $(k+1)!=k!(k+1)$.  This is precisely a
quantifier-free induction invariant after clearing denominators.
\end{proof}

\subsection{A uniform ratio-tail certificate}

Convergence is estimated with the ordinary complex modulus.  A rational
certificate $\abs w\le W$ is checked without a square root by
$(\Re w)^2+(\Im w)^2\le W^2$; center uncertainty is absorbed into $W$ by the
triangle inequality.  The $\ell^1$ norm remains the representation used for
output balls, and $\norm z_1\le2\abs z$ costs only one guard bit.

Let $M_i\ge\abs{A_i}$, and let $L_j(k)>0$ be rational lower bounds for
$\abs{B_j+k}$.  When $t_k\ne0$, the coefficient recurrence yields
\begin{equation}\label{eq:ratio-majorant}
 \frac{\abs{t_{k+1}}}{\abs{t_k}}
 \le W\,
 \frac{\prod_i(k+M_i)}{(k+1)\prod_jL_j(k)}.
\end{equation}
If $t_k=0$, every later term is zero under denominator apartness, so the same
tail conclusion is immediate.  If a rational $0<\rho<1$ and an index $N$ make
the right side at most $\rho$ for every $k\ge N$, then
\begin{equation}\label{eq:geometric-tail}
 \abs{\sum_{k=N}^{\infty}t_k}
 \le\frac{\abs{t_N}}{1-\rho}.
\end{equation}
The proof is finite induction followed by the finite geometric-sum identity;
for a requested precision one truncates the geometric sum after a further
explicit number of terms and bounds its last term.  Converting the final bound
to the $\ell^1$ ball uses the factor $2$ above.  No completed infinite sum is
needed internally.

\begin{theorem}[Primitive-recursive convergence]\label{thm:hyp-conv}
For fixed parameter names with denominator-apartness certificates:
\begin{enumerate}[label=(\alph*)]
\item if $r\le s$, then for every bounded $w$ a primitive-recursive search
      produces $N$ and $\rho=1/2$ satisfying \eqref{eq:ratio-majorant};
\item if $r=s+1$ and a squared-modulus certificate proves $\abs w\le R<1$,
      choose $R<\rho<1$; a primitive-recursive search produces $N$ satisfying
      the ratio bound;
\item if one upper parameter is $-h$ with $h\in\N$, then $t_k=0$ for $k>h$,
      so the series is a finite polynomial without any convergence principle.
\end{enumerate}
The resulting truncation modulus is primitive recursive in the requested
precision and all rational certificate data.
\end{theorem}
\begin{proof}
After multiplying by all positive denominator bounds, the desired ratio
inequality is a polynomial inequality in $k$.  In case (a) the denominator has
strictly larger degree; in case (b) the leading-coefficient ratio is bounded by
$R<\rho$.  A coarse explicit threshold is obtained by replacing every
$k+M_i$ by $(1+M)k$ and each eventually positive $L_j(k)$ by $k/2$.
The finitely many preceding indices are checked directly.  A bounded search up
to the resulting threshold is primitive recursive.  Equation
\eqref{eq:geometric-tail} supplies the modulus.  Case (c) follows from the
factor $-h+h=0$.
\end{proof}

\subsection{Termwise algebra and derivatives}

Uniform ratio bounds on a rational compact set persist after multiplying the
$k$th term by any fixed polynomial in $k$.  Therefore one may differentiate an
externally fixed number of times, multiply series, or shift parameters while
retaining a primitive-recursive tail bound.

\begin{lemma}[Coefficientwise identity principle]\label{lem:coeff-id}
Suppose two certified power or power-log series on the same branch have equal
coefficients by a quantifier-free arithmetic proof and have uniform tail
moduli.  Then their evaluations are ball-equal to every requested precision.
\end{lemma}
\begin{proof}
Given $k$, choose a common truncation index whose two tails are below
$2^{-k-2}$.  The finite partial sums are exactly equal after clearing
 denominators.  Two applications of the triangle inequality give error below
$2^{-k-1}$, which is stronger than required after guard-bit reindexing.
\end{proof}
This finite lemma replaces almost every appeal to the classical identity
theorem in the local part of the report.

\subsection{Terminating and singular parameter cases}

A lower parameter in $\{0,-1,-2,\ldots\}$ makes the naive recurrence undefined.
There are three constructive resolutions, each recorded in the datum:
\begin{enumerate}
\item an upper $-h$ terminates earlier, so all used denominators are nonzero;
\item a common factor is cancelled symbolically before evaluation;
\item a confluent limit is represented by finite parameter jets and a
      quantitative divided-difference bound.
\end{enumerate}
The polynomial formula on the source page is needed because the naive Gamma
prefactor in the one-$G$ representation may vanish exactly when the
hypergeometric series terminates.  We prove the repaired two-$G$ formula in
\cref{sec:hypergeom-conversion} by finite coefficient extraction.
